% STATUS: rewritten 2026-10-08
% Sources read in full: Hochschild 1942 (pp. 678--679, 688--689), Hochschild 1945 (Theorem 4.1), Weibel 1994 (Definition 6.4.1, Example 6.4.2, Lemma 9.2.1), Kharchenko 1991 (pp. 9--10), Humphreys 1972 (Section 5.3), Bresar, Hanselka, Klep and Volcic 2018 (Corollary 6.2), Baclawski 1972 (Theorem 2, p. 354), Benkart, Lopes and Ondrus 2015 (pp. 47--48).
\section{Inner, principal and outer derivations; crossed homomorphisms}\label{sec:basic-theory:inner-outer}

We ask now which derivations come from the algebra itself. Hochschild answers
with a picture from automorphisms \cite[p.~678]{Hochschild1942}. Take an
invertible element $c$. We call the map $a \mapsto cac^{-1}$ an inner automorphism.
Let $c = \exp(ta_0)$, where $t$ is a parameter. Then the automorphism equals
$a + t(a_0a - aa_0)$ plus a term that is small compared with $t$
\cite[pp.~678--679]{Hochschild1942}. Its derivative at $t=0$ is the map
$a \mapsto a_0a - aa_0$.

\begin{definition}[{\cite[Definition~3.2, p.~689]{Hochschild1942}}]\label{def:basic:inner}
Let $A$ be an associative algebra. The derivation $a \mapsto a_0a - aa_0$,
for a fixed $a_0 \in A$, is an \emph{inner derivation} of $A$. For a Lie
algebra $L$, the derivation $l \mapsto [l_0,l]$ is an inner derivation of $L$.
\end{definition}

The next example gives the inner derivations of a ring.

\begin{example}[{\cite[Example~1.3]{GuptaKaurSingh2023}}]\label{ex:basic:inner}
Let $R$ be a ring and $a\in R$. The map $\partial_a(x)=xa-ax$ is a derivation of
$R$, the inner derivation with respect to $a$. It restricts to the zero
derivation on the center of $R$. In particular, if $R$ is commutative, then
$\partial_a$ is zero.
\end{example}

A derivation that is not inner is \emph{outer} \cite[pp.~9--10]{Kharchenko1991}. We state the exact facts for a
ring $R$, following Kharchenko.

The map $b \mapsto \mathrm{ad}\,b$ sends $R$ to $\mathrm{Der}(R)$ and respects
the Lie bracket \cite[pp.~9--10]{Kharchenko1991}. Its kernel is the centre
\cite[pp.~9--10]{Kharchenko1991}. The
image is an ideal, because $[D,\mathrm{ad}\,b] = \mathrm{ad}\,D(b)$ for every
derivation $D$. So the outer part is the quotient $\mathrm{Der}(R)/\mathrm{ad}\,R$.

Humphreys recalls the same fact for a Lie algebra $L$
\cite[\S5.3, p.~23]{Humphreys1972}. The set $\mathrm{ad}\,L$ is an ideal in
$\mathrm{Der}\,L$, because $[\delta,\mathrm{ad}\,x]=\mathrm{ad}(\delta x)$
\cite[\S5.3, p.~23]{Humphreys1972}.

Gupta, Kaur and Singh give a derivation of a matrix algebra that comes from a
derivation of the field \cite[Example~1.4]{GuptaKaurSingh2023}.

\begin{example}[{\cite[Example~1.4]{GuptaKaurSingh2023}}]\label{ex:basic:matrices}
Let $(k,d_k)$ be a differential field. The map $\delta_c$ on $M_m(k)$ defined by
$\delta_c((a_{ij}))=(d_k(a_{ij}))$ is a derivation of $M_m(k)$, the
coordinate-wise derivation.
\end{example}

The same object has other names, and Chapter~\ref{ch:origins} met them. Let a
group $G$ act on an abelian group $A$. Weibel calls a map $D$ with
$D(gh) = gD(h) + D(g)$ a derivation, or crossed homomorphism. The maps
$D_a(g) = ga - a$ are the principal derivations
\cite[Definition~6.4.1, Example~6.4.2]{Weibel1994}. For a ring $R$ and a
bimodule $M$, the principal derivations are the maps $f_m(r) = rm - mr$. In both
settings the group $H^1$ is the quotient of all derivations by the principal
ones \cite[Lemma~9.2.1, p.~301]{Weibel1994}.

So inner, principal and outer derivations are three words for one idea. The
table shows how the terms match.

\begin{center}
\begin{tabular}{lll}
Rings and algebras & Groups & Cohomology \\ \hline
derivation & crossed homomorphism & $1$-cocycle \\
inner derivation & principal derivation & coboundary \\
outer derivation class & element of $H^1$ & cohomology class \\
\end{tabular}
\end{center}

The table gives a second reading of our question. Asking whether every
derivation is inner means asking whether $H^1$ is zero.

Hochschild settled this for algebras over a field
\cite[Theorem~4.1]{Hochschild1945}. His theorem says that an algebra is
separable if and only if $H^1$ vanishes for every two-sided module. So an algebra
is separable exactly when every derivation of it into every bimodule is inner.

We begin with cases in which every derivation is inner. We take first a central
simple algebra, such as the algebra of all $n \times n$ matrices over a field.
Brešar, Hanselka, Klep and Volčič prove the following statement (Theorem~\ref{thm:basic:csa-inner})
\cite[Corollary~6.2]{BresarHanselkaKlepVolcic2018}.\begin{theorem}[{\cite[Corollary~6.2]{BresarHanselkaKlepVolcic2018}}]\label{thm:basic:csa-inner}
Every derivation from a central simple algebra $R$ into an arbitrary $R$-bimodule
$M$ is inner.
\end{theorem}

In particular every derivation of a full matrix algebra over a field is inner. We take second a semisimple Lie algebra. Over an algebraically closed field of characteristic
$0$, such as the complex numbers, every derivation of it is inner
\cite[\S5.3, p.~23]{Humphreys1972}.

Some derivations are not inner, and yet we can list them all. Baclawski does
this for the incidence algebra of an ordered set \cite[p.~351]{Baclawski1972}.
Each derivation is the sum of an inner derivation and a second kind, which comes
from an additive function \cite[Theorem~2, p.~353]{Baclawski1972}. When the
ordered set has a smallest or a largest element, every derivation is inner
\cite[p.~354]{Baclawski1972}.

We see next that the answer can also depend on the field. Consider the Weyl
algebra of differential operators with polynomial coefficients
\cite[p.~46]{BenkartLopesOndrus2015}. Over a field of characteristic $0$ all its
derivations are inner \cite[p.~48]{BenkartLopesOndrus2015}. Over a field of
prime characteristic some are not \cite[p.~48]{BenkartLopesOndrus2015}.

We return to $H^1$ in Chapter~\ref{ch:cohomology}. So inner derivations are the
derivatives of inner automorphisms.
Section~\ref{sec:basic-theory:exponential} follows this link between
derivations and automorphisms in general.
